\documentclass[12pt,a4paper]{article}
\usepackage[utf8]{inputenc}
\usepackage[indonesian]{babel}
\usepackage{amsmath}
\usepackage{amsfonts}
\usepackage{amssymb}
\usepackage{geometry}

\geometry{margin=2.5cm}

\title{Menghitung Volume Benda Putar dengan Teorema Pappus}
\author{Ahmad Fakhrul Bawani}
\date{}

\begin{document}

\maketitle

\section*{Soal}

Use a theorem of Pappus to find the volume generated by revolving about the line $x = 7$ the triangular region bounded by the coordinate axes and the line $5x + y = 30$.

\subsection*{1. Rumus Umum Teorema Pappus}
Teorema Pappus untuk volume menyatakan bahwa jika sebuah daerah bidang datar dengan luas $A$ diputar mengelilingi sebuah poros yang tidak memotong daerah tersebut, maka volume benda putar ($V$) yang dihasilkan adalah:
\begin{equation*}
  V = 2\pi \cdot R \cdot A
\end{equation*}
Di mana $R$ adalah jarak tegak lurus dari pusat massa (sentroid) daerah tersebut ke garis poros putaran.

---

\subsection*{2. Menentukan Karakteristik Daerah Segitiga}
Garis pembatas daerah adalah $5x + y = 30$. Kita tentukan titik potong segitiga tersebut dengan sumbu-sumbu koordinat:
\begin{itemize}
  \item \textbf{Titik potong sumbu-}x ($y = 0$): $5x = 30 \implies x = 6$. Maka koordinat alasnya di $(6,0)$.
  \item \textbf{Titik potong sumbu-}y ($x = 0$): $y = 30$. Maka koordinat tingginya di $(0,30)$.
\end{itemize}
Segitiga yang terbentuk dibatasi oleh titik-titik sudut $(0,0)$, $(6,0)$, dan $(0,30)$.

\begin{itemize}
  \item \textbf{Luas Segitiga ($A$):}
    \begin{equation*}
      A = \frac{1}{2} \cdot \text{alas} \cdot \text{tinggi} = \frac{1}{2} \cdot 6 \cdot 30 = 90
    \end{equation*}

  \item \textbf{Menentukan Sentroid Segitiga ($\bar{x}, \bar{y}$):}
    Berdasarkan catatan soal, koordinat sentroid sebuah segitiga dengan titik-titik sudut $(x_1, y_1)$, $(x_2, y_2)$, dan $(x_3, y_3)$ berada tepat di rata-rata titik sudutnya:
    \begin{equation*}
      \bar{x} = \frac{x_1 + x_2 + x_3}{3} = \frac{0 + 6 + 0}{3} = 2
    \end{equation*}
    *(Kita hanya membutuhkan nilai $\bar{x}$ karena daerah diputar terhadap garis vertikal).*
\end{itemize}

---

\subsection*{3. Jarak Putar ($R$) dan Perhitungan Volume}
Daerah diputar mengelilingi garis vertikal $x = 7$. Karena sentroid daerah berada di $\bar{x} = 2$, jarak horizontal $R$ dari sentroid ke sumbu putar adalah:
\begin{equation*}
  R = 7 - \bar{x} = 7 - 2 = 5
\end{equation*}
Perhatikan bahwa garis $x = 7$ tidak memotong segitiga (karena batas segitiga terjauh hanya sampai $x = 6$), sehingga syarat Teorema Pappus terpenuhi.

Sekarang, kita hitung volume benda putar ($V$):
\begin{align*}
  V &= 2\pi \cdot R \cdot A \\
  V &= 2\pi \cdot 5 \cdot 90 \\
  V &= 900\pi
\end{align*}

\subsection*{Kesimpulan}
Jadi, volume benda putar yang dihasilkan adalah \textbf{$900\pi$ satuan volume}.

\end{document}